\section{Một khảo sát thấy chỗ hổng và gọi sai tên nó}
\label{sec:ch17-khao-sat-biet-cho-hong}

\index{chuỗi suy luận!khảo sát 2026 đặt độ trung thực dưới mặt nước|(}%
Mục trước để lại một khảo sát không thấy chỗ hổng. Mục này đọc một khảo sát thấy
nó, nói ra ba lần, và vẽ hẳn một hình cho nó. Bài này
chương~\ref{ch:mo-hinh-can-giai-thich} đã đọc mục 3 của nó để mượn cách nhìn
rằng suy luận là hành vi đo được dưới một điều kiện; ở đây nó được mở lại cho
một câu hỏi khác, và câu hỏi ấy được trả lời ở trang 76 của 103 trang.

\index{chuỗi suy luận!ba lần khảo sát nêu chỗ nghi}%
Ba chỗ nêu chỗ nghi nằm cách xa nhau và nói cùng một điều. Phần mở đầu, trang 2,
viết rằng tuy các kỹ thuật như nhắc theo chuỗi suy luận phơi ra được các dấu vết
suy luận trung gian, vẫn còn đáng nghi ngờ liệu những dấu vết ấy có phản ánh
trung thực các quá trình tính toán nằm dưới hay chỉ là những lời biện minh hậu
nghiệm nghe hợp lý~\autocite{p08periodic}. Mục về chuỗi suy luận, trang 8, viết
rằng cách nhắc ấy không bảo đảm tính đúng đắn hay độ trung thực của quá trình
suy luận sinh ra, vì mô hình vẫn vận hành bằng phép sinh token kế tiếp theo xác
suất~\autocite{p08periodic}. Và danh sách hướng dẫn thực hành ở trang 15 khép
lại bằng câu rằng chuỗi dài hơn có thể cải thiện kết quả, nhưng riêng việc dài
lời thì không bảo đảm suy luận trung thực hay đúng~\autocite{p08periodic}.

Câu đi ngay sau câu thứ hai đáng đọc kỹ, vì nó là chỗ khảo sát này giống khảo
sát của mục trước. Sau khi nêu chỗ nghi, bài viết tiếp rằng dù vậy, cấu trúc
theo từng bước vẫn cải thiện khả năng diễn giải và cho cơ hội kiểm tra cục bộ và
phân tích lỗi, so với cách nhắc trả lời thẳng~\autocite{p08periodic}. Hai câu
ấy đứng trong cùng một đoạn. Câu trước nói không ai biết dấu vết có phản ánh
tính toán hay không; câu sau nói dấu vết cải thiện khả năng diễn giải; và bài
không nói khả năng diễn giải ấy đứng trên cái gì sau khi câu trước đã được chấp
nhận.

\index{chuỗi suy luận!hình mặt nước}%
Hình 27, trang 76, là chỗ bài nói rõ nhất. Nó vẽ một tảng băng: bên trên mặt
nước là độ chính xác, dán nhãn là thứ mà các bộ đánh giá đo; bên dưới mặt nước
là sáu tính chất, và tính chất đầu tiên trong sáu là độ trung
thực~\autocite{p08periodic}. Bài gọi tên cả dụng cụ: dưới mục ấy, hình liệt kê
ba họ, can thiệp vào dấu vết, chuỗi suy luận \tn{phản thực}{counterfactual}, và
\tn{phân tích trung gian nhân quả}{causal mediation analysis}. Khung hướng dẫn ở chân hình thì viết rằng nếu một bài chỉ báo cáo độ
chính xác thì hãy coi mọi tuyên bố của nó về \enquote{suy luận} là sơ
bộ~\autocite{p08periodic}.

Ba họ dụng cụ ấy đáng đặt cạnh chương~\ref{ch:chi-so-khong-do-do-trung-thuc}, vì
cả ba đều đã được đem ra kiểm ở đó. Bốn trong tám chỉ số của
mục~\ref{sec:ch09-ket-qua}, gồm Adding Mistakes, Early Answering, Filler Tokens
và SCM, đều can thiệp vào chuỗi rồi xem đáp án có đổi không, tức là hai họ đầu mà
hình 27 gọi tên; và SCM thì áp đúng phép phân tích trung gian nhân quả, tức họ
thứ ba. Bảng~\ref{tab:ch09-auroc} in điểm của cả bốn. Khảo sát 2026 gọi tên ba họ
dụng cụ, không nói họ nào trong ba đã được kiểm, và cả ba đều đã được kiểm rồi.

\index{mức quan trọng!khảo sát 2026 định nghĩa nhầm vào đây}%
Nhưng chỗ đáng ghi nhất của hình 27 là câu định nghĩa. Bên trong hình, độ trung
thực được chú bằng đúng một câu hỏi, dấu vết có gây ra đáp án hay không; và ở
bảng chú giải bên phải, câu ấy được viết đầy đủ thành: chuỗi suy luận có thật sự
gây ra đáp án, hay nó là đồ trang trí hậu nghiệm~\autocite{p08periodic}.

Câu hỏi ấy không phải câu hỏi của định nghĩa~\ref{def:ch01-do-trung-thuc}. Nó là
câu hỏi về \tn{mức quan trọng}{importance}, và
mục~\ref{sec:ch09-dinh-nghia-lai} đã tách hai tính chất ấy ra bằng hai ví dụ
chạy theo hai chiều ngược nhau, một bước trung thực mà không quan trọng và một
bước quan trọng mà không trung thực, rồi ghi rằng đồng nhất chúng sẽ đảo đúng cái
ranh giới mà việc đánh giá độ trung thực muốn nắm.

\index{chuỗi suy luận!khảo sát 2026 đặt độ trung thực dưới mặt nước|)}%
Vậy một khảo sát năm 2026, trong đúng cái hình dựng ra để nói rằng các bộ đánh
giá đang đo nhầm thứ, định nghĩa độ trung thực bằng chính tính chất mà chẩn đoán
trung tâm của bài neo bảo là hay bị lẫn với nó. Sách đã đặt tên riêng cho hai
tính chất ấy từ chương~\ref{ch:chi-so-khong-do-do-trung-thuc}, và lý do đặt tên
là để không lẫn. Đây là lần đầu sách gặp phép lẫn ấy nằm ngay trong một định
nghĩa được in ra, chứ không nằm bên trong một chỉ số.
